% Modernised reading — fonds Alexandre Grothendieck, shelfmark 44, pages 1-8.
%
% This is an INTERPRETATION, not a transcription. Everything the critical
% apparatus carried moves into footnotes. In any dispute of substance the
% transcription governs: whoever cites Grothendieck must cite batch-01.fr.tex.
%
% Pass: Opus 5 (claude-opus-5), 2026-08-30 — first pass, unchecked against the
% pages by a human. The folder was read whole; it holds one batch, pages 1-8.
%
% Notation fixed for the whole document. The folder writes the intermediate
% Jacobian two ways: P^i(X) on the three sheets of the programme (pages 1, 3,
% 5) and J^i(X) in the letter (pages 7, 8). J^i(X) is used throughout here and
% the variance is footnoted once. Z^i(X) is the group of codimension-i cycles,
% Z^i_alg those algebraically equivalent to zero, Z^i_alb those the folder
% calls "albanesiennement" equivalent to zero; the transcription notes that his
% Z is drawn now angular now curly without distinction, and that the subscript
% alb is circled at both its occurrences.
%
% Substantive departures from the pages, each footnoted where it occurs:
% — Th. 4 (page 1) is stated on the page without the word "primitive", and he
%   cancels it: "Cancelé tel quel". The margin of page 7 says what to put in
%   its place, and it is the classical restriction to primitive cohomology.
%   The reading states the corrected form and footnotes the page's.
% — Th. 6 (page 3) has its hypothesis 2i <= n+1 struck through and never
%   replaced. The range under which the statement is true is supplied and
%   footnoted.
% — Th. 7 (page 5) asserts numerical equivalence = tau-equivalence for
%   algebraic cycles without restriction on codimension. In that generality it
%   is false — Griffiths (1969). The reading states the divisor case, which is
%   Matsusaka's theorem and is what the marginal name records, and footnotes
%   the rest.
% — Th. 5 (page 3) is the most heavily marked statement of the folder; the
%   group in which its pairing takes values is not recoverable from the page,
%   and the reading says so rather than supplying one.
%
% What the folder announces and never establishes: the whole of it. These are
% seven statements offered as a programme, with no proof on any leaf, and the
% letter says as much — "je n'ai pas prouvé grand-chose dessus pour l'instant".
% Pages 2 and 4 belong to two other documents of his own and are read here as
% an appendix, in the terms of their own subjects.
\documentclass[11pt,a4paper]{article}
% Le préambule commun à toutes les transcriptions du fonds Grothendieck.
%
% Il tient en une idée : **l'incertitude se compose, elle ne se résout pas.**
% Une transcription qui lisse un mot douteux a détruit la seule chose pour
% laquelle on la faisait — un lecteur doit pouvoir savoir, à chaque endroit,
% ce qui était sur le feuillet et ce qui a été ajouté. Les six macros qui
% suivent sont donc l'appareil critique, et non de la décoration.
%
% Les mêmes commandes sont reconnues par `scripts/render.mjs`, qui produit la
% vue de lecture du site. Une macro ajoutée ici doit l'être aussi là-bas,
% faute de quoi le rendu échouera bruyamment — ce qui est voulu : un rendu
% silencieusement faux est pire qu'un rendu absent.

\ProvidesPackage{grothendieck}[2026/08/08 Transcriptions du fonds Grothendieck]

\usepackage{amsmath,amssymb,amsthm}
\usepackage{mathrsfs}
\usepackage{stmaryrd}
% Les diagrammes commutatifs sont partout dans ce fonds : c'est la langue
% même de la géométrie algébrique de Grothendieck, et une transcription qui
% les rendrait en prose ne transcrirait rien.
\usepackage{tikz-cd}
% Français par défaut — c'est la langue des feuillets — mais l'anglais est
% chargé aussi : la traduction et le résumé sont des documents anglais, et la
% typographie française (espaces avant les deux-points) y serait une faute.
% Ces documents basculent par \selectlanguage{english} en tête de corps.
\usepackage[english,french]{babel}
\usepackage{fontspec}
\usepackage{geometry}
\geometry{a4paper,margin=25mm,marginparwidth=28mm}
\usepackage{marginnote}
\usepackage{xcolor}
% ulem, not soul. `soul`'s \st and \ul are built on a character-by-character
% scan that XeLaTeX's font machinery breaks: the document fails with an
% undefined control sequence rather than degrading. ulem does the same two jobs
% and is Unicode-engine safe. `normalem` stops it hijacking \emph, which the
% transcriptions use for Grothendieck's own emphasis.
\usepackage[normalem]{ulem}
\usepackage{hyperref}
\hypersetup{colorlinks=true,linkcolor=black,urlcolor=black,pdfborder={0 0 0}}

% --- Métadonnées du lot -----------------------------------------------------
% Reprises telles quelles par le script de rendu, qui les affiche en tête de
% la vue de lecture. Elles fixent ce que le fichier prétend couvrir : sans
% elles, une transcription flotte sans point d'attache dans un fonds de seize
% mille feuillets.

\newcommand{\folder}[1]{\gdef\grothFolder{#1}}
\newcommand{\batch}[1]{\gdef\grothBatch{#1}}
\newcommand{\leaves}[2]{\gdef\grothFirst{#1}\gdef\grothLast{#2}}
\newcommand{\foldertitle}[1]{\gdef\grothFoldertitle{#1}}
\gdef\grothFolder{?}\gdef\grothBatch{?}\gdef\grothFirst{?}\gdef\grothLast{?}\gdef\grothFoldertitle{}

% --- Le résumé --------------------------------------------------------------
%
% Il ouvre la lecture modernisée, sous le titre « Résumé », et il est composé
% en petit. Ce n'est pas un ornement : c'est le seul passage du document qui
% ne soit pas des mathématiques, et un lecteur doit pouvoir le reconnaître
% avant de l'avoir lu — puis le sauter s'il connaît déjà le sujet, ou s'y
% arrêter s'il ne le connaît pas. Le filet qui le suit marque la reprise du
% corps.

\newenvironment{resume}
  {\small\setlength{\parskip}{0.45em}}
  {\par\medskip\hrule\medskip}

% --- Mots-clés ---------------------------------------------------------------
%
% En anglais, en clôture du résumé de la lecture modernisée : le vocabulaire
% moderne sous lequel chercher ce que les feuillets construisent. Le manifeste
% du site (`npm run manifest`) les extrait pour étiqueter la cote entière —
% cette ligne est la source unique des tags, il n'y a pas de fichier à côté.
\newcommand{\keywords}[1]{\par\smallskip\noindent
  {\footnotesize\textsc{Keywords} --- \textit{#1}}\par}

% --- L'appareil critique ----------------------------------------------------

% Le numéro de feuillet, en marge. C'est l'ancre qui permet de régler un
% désaccord sur une formule en regardant *un* feuillet plutôt qu'une chemise
% de six cents. Le site s'en sert aussi pour tourner les pages du fac-similé
% quand on fait défiler la transcription.
\newcommand{\leaf}[1]{%
  \marginnote{\footnotesize\color{gray!70}\textsf{#1}}%
  \label{leaf:#1}\ignorespaces}

% Illisible : on ne devine pas. Un mot inventé qui se lit comme les autres est
% le pire résultat possible.
\newcommand{\ill}{\textcolor{red!60!black}{[\,\ldots\,]}}

% Lecture douteuse : le mot est donné, et signalé comme incertain.
\newcommand{\uncertain}[1]{\uline{#1}}

% Ajout de l'éditeur — une lettre manquante, un mot sous-entendu.
\newcommand{\add}[1]{\textcolor{blue!50!black}{[#1]}}

% Biffé par Grothendieck lui-même. Ce qu'il a rayé fait partie du document :
% une rature dit souvent où la pensée a changé de direction.
\newcommand{\struck}[1]{\sout{#1}}

% Note du transcripteur, sur l'état matériel du feuillet.
\newcommand{\note}[1]{\par\smallskip{\small\color{gray!60!black}\textit{[#1]}}\par\smallskip}

% Note marginale de Grothendieck, distincte du corps du texte.
\newcommand{\marginal}[1]{\par\smallskip\noindent\hspace{1em}%
  {\small\color{gray!60!black}#1}\par\smallskip}

% --- Titre ------------------------------------------------------------------

\AtBeginDocument{%
  \begin{center}
    {\large\bfseries Fonds Alexandre Grothendieck --- cote n\textsuperscript{o}~\grothFolder}\\[2pt]
    {\itshape\grothFoldertitle}\\[4pt]
    {\small feuillets \grothFirst--\grothLast\ (lot \grothBatch)}\\[2pt]
    {\footnotesize\color{gray!60!black}Transcription établie d'après le fac-similé numérisé
     par l'Université de Montpellier.\\
     Les lectures incertaines sont soulignées, les illisibles notées [\,\ldots\,],
     les ajouts entre crochets.}
  \end{center}
  \vspace{1em}}

\endinput


\folder{44}
\pages{1}{8}
\dating{1964}
\watermark{Édition de démonstration\\interprétation personnelle de l'œuvre}
\foldertitle{Jacobiennes intermédiaires : notes manuscrites (s.d.), lettre (1964) — lecture modernisée du dossier entier}

\begin{document}

\section*{Résumé}

\begin{resume}
Une courbe porte sa jacobienne. Ces huit feuillets demandent ce qu'une variété
de dimension quelconque porte à sa place, et combien.

À une courbe algébrique on sait associer, depuis le
\textsc{xix}\textsuperscript{e} siècle, un objet géométrique dont les points
s'additionnent — sa \emph{jacobienne} — et qui classe ses systèmes de points :
deux configurations de points de la courbe donnent le même point de la
jacobienne exactement quand elles sont reliées par une famille continue. En
dimension supérieure, deux cas restent faciles, et ce sont les deux extrêmes.
Les diviseurs — les sous-objets de codimension $1$, ceux qu'une seule équation
découpe — sont classés par la \emph{variété de Picard}. Les points — la
codimension maximale — sont classés par la \emph{variété d'Albanese}. Entre les
deux, pour les sous-objets de codimension $2$, $3$, \ldots, on ne disposait de
rien d'algébrique.

Sur les nombres complexes on disposait tout de même de quelque chose. Weil
avait construit, par des moyens transcendants — des intégrales, des périodes,
la topologie de l'espace sous-jacent —, des tores qu'on appelle
\emph{jacobiennes intermédiaires} : un par codimension. La construction est
belle et elle ne sert à rien hors des complexes, parce qu'elle passe par des
objets qui n'existent pas ailleurs. Or c'est ailleurs que Grothendieck
travaille : sur un corps quelconque, où l'on ne peut ni intégrer ni parler de
petits lacets.

L'idée de ces feuillets est de renverser la définition. Plutôt que de
construire l'objet et de constater ensuite qu'il classe les cycles, on part de
ce qu'il doit faire. Une famille de cycles paramétrée par une variété $S$
donne, par la construction d'Albanese appliquée à $S$, un morphisme d'une
variété abélienne vers le groupe des cycles ; le premier théorème dit que le
noyau de ce morphisme est exactement ce qu'on croit, et le deuxième que
l'image, quelle que soit la famille, ne dépasse jamais une taille fixée à
l'avance par la topologie de $X$. Les deux ensemble donnent l'existence d'une
plus grande image possible : c'est elle, $J^i(X)$, la jacobienne intermédiaire
algébrique. La relation d'équivalence entre cycles n'est plus choisie d'abord ;
elle est \emph{ce que cet objet universel ne distingue pas}, et Grothendieck la
baptise « équivalence d'Albanese » sans savoir la décrire autrement.

Restent les propriétés qu'on veut, et qui font les cinq autres théorèmes : que
la jacobienne de codimension $i$ et celle de codimension $n+1-i$ soient duales
l'une de l'autre ; qu'au milieu, quand $i$ et $n+1-i$ coïncident, cette
autodualité provienne d'une \emph{polarisation}, ce qui est la forme
algébrique du théorème de positivité de Hodge ; qu'une section hyperplane
générique perde peu d'information ; et que l'équivalence numérique — la plus
grossière possible, celle qui ne retient d'un cycle que ses nombres
d'intersection — coïncide avec l'équivalence algébrique à torsion près.

La lettre à Serre du 12 août 1964, qui occupe les deux derniers feuillets et
que Grothendieck a numérotée à la suite du programme, isole le seul de ces
points qu'il ne sait pas franchir, et le réduit à un énoncé « bien terre à
terre » sur les variétés abéliennes : une certaine correspondance symétrique
est-elle positive ? Il demande si les spécialistes sauraient répondre. C'est
aussi le feuillet où on le voit travailler : « je suis moralement sûr que ce
sont les bonnes, bien que je n'ai pas prouvé grand-chose dessus pour
l'instant ; par contre, j'ai des conjectures en pagaille ». La certitude
précède la démonstration, et il le dit sans embarras.

Ce que le programme demandait a été en partie obtenu et en partie réfuté. La
jacobienne intermédiaire algébrique existe : on l'appelle aujourd'hui le
\emph{représentant algébrique} du groupe des cycles, et c'est exactement l'objet
que les deux premiers théorèmes appellent. Le dernier théorème, en revanche,
est faux dès la codimension $2$ : cinq ans après cette lettre, Griffiths
exhibait des courbes sur une hypersurface qui sont numériquement triviales sans
qu'aucun multiple d'elles ne soit algébriquement trivial. L'écart entre les
deux équivalences, que ce feuillet croyait nul, porte aujourd'hui le nom de
groupe de Griffiths.

\keywords{intermediate Jacobian, algebraic representative, regular
homomorphism, Albanese equivalence, algebraic equivalence, numerical
equivalence, tau-equivalence, Griffiths group, Abel-Jacobi map, Hodge-Riemann
bilinear relations, primitive cohomology, polarisation, Neron-Severi group,
divisorial correspondence, Poincare biextension, weak Lefschetz theorem, Tate
module, Betti number, Matsusaka theorem, graded ring completion, root system,
root datum, Cartan integers}
\end{resume}

\section*{Le fil du dossier, et l'ordre des feuillets}

Le dossier est une seule liasse, paginée de $1$ à $5$ de sa main sur les
feuillets qu'il a écrits. Trois d'entre eux — les pages 1, 3 et 5 de
l'archiviste — portent un « Programme de travail » en sept théorèmes numérotés.
Les deux derniers — pages 7 et 8 — portent le tapuscrit de sa lettre à Serre du
12 août 1964, qu'il a numérotée $4$ et $5$ pour qu'elle ferme la même suite. Le
programme et la lettre disent la même chose deux fois, une fois en énoncés secs
et une fois en prose, et la lettre est la plus explicite des deux.

La colonne vertébrale est celle-ci. Les théorèmes 1 et 2 \emph{construisent}
l'objet : le premier identifie ce qu'une famille de cycles apporte, le second
borne uniformément ce que toutes les familles peuvent apporter, et cette borne
uniforme est ce qui permet de prendre la plus grande. Les théorèmes 3 et 4 lui
demandent la structure qu'on veut : une dualité entre codimensions
complémentaires, et au milieu une polarisation. Les théorèmes 5 et 6 le
comparent à ce qu'on obtient sur une section hyperplane. Le théorème 7 est d'un
autre ordre : il ne parle plus de la jacobienne mais de la relation
d'équivalence la plus grossière, et c'est le seul des sept dont on sait
aujourd'hui qu'il est faux.

Deux feuillets n'appartiennent pas à cette liasse et n'ont rien à voir avec
elle. Il a écrit son programme au dos de papier déjà servi : la page 2 est un
feuillet à lui, barré, sur la complétion d'un anneau gradué, et la page 4 un
feuillet à lui, barré lui aussi, sur les axiomes d'une donnée radicielle. Ils
sont lus ici en appendice, dans les termes de leurs propres sujets, et rien
n'est affirmé d'une continuité entre eux et les jacobiennes.

\textbf{Notation.} Les trois feuillets du programme écrivent $P^{i}(X)$ ce que
la lettre écrit $J^{i}(X)$ : c'est le même objet, et $J^{i}(X)$ est employé
partout ici.\footnote{Le changement de lettre est sur les pages ; il n'est
expliqué nulle part. Le transcripteur note par ailleurs que la lettre des
groupes de cycles est tracée tantôt en $Z$ anguleux tantôt en $\mathcal{Z}$
bouclé, sans distinction apparente, et que l'indice \emph{alb} est entouré d'un
cercle aux deux endroits où il paraît — sans doute la marque d'une notation
qu'il tenait encore pour provisoire.} $X$ est partout une variété projective
lisse connexe de dimension $n$ sur un corps algébriquement clos $k$ ;
$\mathcal{Z}^{i}(X)$ le groupe des cycles de codimension $i$,
$\mathcal{Z}^{i}_{\mathrm{alg}}(X)$ ceux qui sont algébriquement équivalents à
zéro, $\mathcal{Z}^{i}_{\mathrm{alb}}(X)$ ceux que le dossier dit
« albanesiennement » équivalents à zéro. $b_{j}(X)$ est le $j$-ième nombre de
Betti.

\pagerange{1}{1}
\section*{Ce que la jacobienne intermédiaire doit être (page 1)}

\subsection*{Une famille de cycles, et ce qu'elle engendre}

Soit $S$ une variété de paramètres pointée et $\xi$ une famille de cycles de
codimension $i$ de $X$ paramétrée par $S$ — concrètement, un cycle de
codimension $i$ sur $S\times X$, dont la fibre au point marqué est nulle. La
famille définit une application de $S$ dans le groupe des classes de cycles, et
comme le but est un groupe commutatif, cette application se factorise par la
variété d'Albanese de $S$ : c'est la propriété universelle d'Albanese, et c'est
tout ce qu'on lui demande ici. On obtient

\[
  \varphi_{\xi} : \mathrm{Alb}(S) \longrightarrow
  \mathcal{Z}^{i}(X)/\mathcal{Z}^{i}_{\mathrm{alb}}(X).
\]

Le premier théorème dit que le noyau $N(\xi)$ de ce morphisme est celui qu'on
attend, et l'énoncé prend soin de préciser \emph{ensemblistement} : c'est une
égalité de sous-ensembles de points, pas encore de sous-schémas.\footnote{La
précision « (ensemblistement) » est ajoutée en interligne sur la page, un mot
biffé devant elle. Elle n'est pas anodine en 1964 : la structure schématique
des noyaux est précisément ce qui distingue le point de vue de Grothendieck de
celui de ses prédécesseurs, et il signale ici qu'il ne l'affirme pas.} Le
quotient $\mathrm{Alb}(S)/N(\xi)$ est donc une variété abélienne, notée
$P_{\xi}$ : c'est la part de la jacobienne cherchée que la seule famille $\xi$
donne à voir.

\subsection*{La borne uniforme, et l'existence}

Le second théorème est celui qui fait tout le travail :

\[
  \dim P_{\xi} \;\leq\; \tfrac{1}{2}\, b_{2i-1}(X),
\]

et la borne ne dépend pas de $\xi$. C'est ce qui permet de faire la somme de
toutes les familles à la fois : une famille croissante de sous-variétés
abéliennes d'un groupe, uniformément bornée en dimension, est stationnaire, et
il existe donc un plus grand $P_{\xi}$. Cet objet maximal est la jacobienne
intermédiaire $J^{i}(X)$, et l'équivalence d'Albanese est, par construction, ce
qu'il ne sépare pas.\footnote{C'est exactement l'architecture de la théorie
telle qu'elle s'est faite. Ce que le dossier appelle $J^{i}(X)$ s'appelle
aujourd'hui le \emph{représentant algébrique} du groupe
$A^{i}(X) = \mathcal{Z}^{i}_{\mathrm{alg}}(X)$ modulo équivalence rationnelle :
la variété abélienne universelle pour les \emph{homomorphismes réguliers} de
$A^{i}(X)$ vers les variétés abéliennes, au sens de Samuel. Son existence en
codimension quelconque a été établie par H. Saito (1979) ; le cas de la
codimension $2$, plus fin, est dû à Murre (1983). La démonstration procède par
la borne uniforme, c'est-à-dire par le théorème 2 de cette page.}

La raison de la borne est donnée en une ligne, dans la lettre : le module de
Tate de $J^{i}(X)$ se plonge, à un quotient près, dans la cohomologie de degré
$2i-1$. Nous y venons.

\subsection*{Dualité, et positivité au milieu}

Le troisième théorème demande que le quotient

\[
  \mathcal{Z}^{i}_{\mathrm{alg}}(X)/\mathcal{Z}^{i}_{\mathrm{alb}}(X)
  \;=\; J^{i}(X)
\]

porte une structure de variété abélienne, et que $J^{i}(X)$ et $J^{n+1-i}(X)$
soient canoniquement duales l'une de l'autre. Les deux extrémités de cette
famille d'énoncés sont classiques : pour $i=1$ on a $J^{1}(X)=\mathrm{Pic}^{0}(X)$
et pour $i=n$ on a $J^{n}(X)=\mathrm{Alb}(X)$, et la dualité entre la variété de
Picard et la variété d'Albanese est un théorème.\footnote{La lettre le dit
explicitement, page 7 : « $J^{1} = \mathrm{Pic}^{0}$ et $J^{n} =
\mathrm{Alb}^{0}$ ». Les indices $2i-1$ et $2(n+1-i)-1 = 2n-2i+1$ sont
complémentaires dans $H^{*}(X)$, ce qui est la raison cohomologique de la
dualité demandée.} Ce que demande le théorème 3 est la forme générale de ce
fait.

Le quatrième théorème est le point difficile, et Grothendieck le sait : il
l'annule d'un mot en marge, entouré et fléché vers l'énoncé — « Cancelé tel
quel ». Quand $n = 2m-1$ est impair, la dualité du théorème 3 devient une
\emph{auto}dualité de $J^{m}(X)$, laquelle équivaut à la donnée d'une classe de
correspondance divisorielle symétrique sur $J^{m}(X)\times J^{m}(X)$,
c'est-à-dire d'un élément du groupe de Néron--Severi de $J^{m}(X)$. La question
est de savoir si cet élément est ample, autrement dit si l'autodualité provient
d'une \emph{polarisation}. Sous cette forme l'énoncé est faux, et la marge de la
page 7 dit pourquoi : il faut prendre la classe \emph{primitive}.\footnote{La
note marginale de la page 7 est en partie illisible ; ce qui s'en lit est
« Cancelé tel quel, il faut \ldots\ de la \ldots\ primitive ». La correction
est celle qu'on attend : les relations bilinéaires de Hodge--Riemann ne donnent
la définitude que sur la cohomologie \emph{primitive}, le complément de
l'opérateur de Lefschetz. Sur toute la cohomologie de degré $2m-1$ la forme
d'intersection n'est pas définie, et l'énoncé du feuillet, qui ne dit pas
« primitive », affirme donc plus qu'il n'est vrai. Il l'a vu et l'a écrit
lui-même deux feuillets plus loin.} Corrigé ainsi, l'énoncé est la forme
algébrique du théorème de positivité de Hodge, et c'est lui que la lettre
réduira à une question sur les variétés abéliennes.

Le feuillet nomme en marge, entre guillemets, la provenance de chaque énoncé :
« Hodge » contre le théorème 4, « Néron » et « Lefschetz » contre les théorèmes
5 et 6, « Matsusaka » contre le théorème 7.\footnote{« Hodge » et « Lefschetz »
sont barrés d'un long trait sur les pages, sans qu'on sache si la rature vise
le nom ou l'énoncé qu'il désigne ; dans le cas de « Hodge » elle accompagne le
« Cancelé tel quel ». La lecture de « Matsusaka » n'est pas assurée, mais le
nom reparaît en toutes lettres, dactylographié, à la page 8.}

\pagerange{3}{3}
\section*{Ce qu'une section hyperplane conserve (page 3)}

\subsection*{L'accouplement, et les nombres d'intersection}

Le cinquième théorème compare l'intersection dans $X$ à la dualité dans une
section hyperplane. Soit $Y \subset X$ une section hyperplane générique, de
dimension $n-1$, et soient $Z$, $Z'$ deux cycles de $X$ de codimensions $i$ et
$j$ avec

\[
  i + j \;=\; n \;=\; \dim Y + 1 .
\]

Les restrictions $z = Z\cdot Y$ et $z' = Z'\cdot Y$ sont supposées
algébriquement équivalentes à zéro sur $Y$. Elles définissent alors des points
de $J^{i}(Y)$ et de $J^{j}(Y)$, et la condition $i+j = \dim Y + 1$ est
exactement celle sous laquelle ces deux variétés abéliennes sont duales l'une
de l'autre par le théorème 3. Il y a donc entre elles un accouplement canonique
— celui que le feuillet attribue au « foncteur de Néron--Tate » — et le
théorème affirme qu'il calcule le nombre d'intersection :

\[
  \langle z, z'\rangle \;=\; (-1)^{m}\, Z\cdot Z' .
\]

L'énoncé est le plus raturé du dossier, et deux choses n'en sont pas
récupérables : le groupe dans lequel l'accouplement prend ses valeurs, et
l'entier $m$ qui gouverne le signe.\footnote{La transcription porte ici une
douzaine de marques : la condition sur les codimensions est écrite, biffée,
réécrite au-dessus ; l'expression qui devrait décrire l'accouplement est
illisible ; et $m$ n'est défini nulle part sur le feuillet. La lecture ne
supplée ni l'un ni l'autre. Le cas $n=2$, où $Y$ est une courbe, $i=j=1$ et
$J^{1}(Y)$ la jacobienne de $Y$, est classique et fixe le contenu de l'énoncé :
l'accouplement de dualité entre la jacobienne et sa duale y calcule, au signe
près, l'intersection des deux diviseurs sur la surface. Ce qui, en toute
codimension, joue le rôle de cet accouplement est la biextension de Poincaré ;
la théorie des hauteurs de Néron et Tate, à laquelle renvoie le nom porté sur
le feuillet, en est le versant arithmétique.} Une note marginale ajoute la
vérification de cohérence attendue : lorsque $2i = 2j$, les deux membres sont
nuls.\footnote{La fin de cette note est illisible. La raison est claire : si
$i=j$ alors $n=2i$ est pair, l'accouplement au degré médian est alterné, et un
élément y est orthogonal à lui-même.}

Une seconde note, tracée dans un triangle au bas du feuillet et en grande
partie illisible, dit que l'énoncé devrait résulter d'un invariant plus fin
admettant un « variant local » ; la lecture ne la reconstitue pas.

\subsection*{Un théorème de type Lefschetz}

Le sixième théorème demande, sous les mêmes hypothèses, que la restriction

\[
  J^{i}(X) \longrightarrow J^{i}(Y)
\]

soit une isogénie tant que $2i \leq n-1$, et qu'au degré médian $2i = n$ elle
reste injective à noyau fini, son image étant tout ce que $J^{i}(Y)$ retient de
$X$.\footnote{Sur la page, la condition portée d'abord est $2i \leq n+1$, puis
barrée, et rien ne la remplace : les trois lignes qui suivent sont un
enchevêtrement de variantes toutes biffées sauf une. La borne $2i\leq n-1$ est
restituée ici parce que c'est celle sous laquelle l'énoncé est vrai et celle que
le théorème de Lefschetz faible impose : la restriction
$H^{2i-1}(X)\to H^{2i-1}(Y)$ est un isomorphisme pour $2i-1 \leq n-2$. Le mot
\emph{isogénie}, lui, est souligné sur la page et n'est pas douteux ; c'est le
seul mot du feuillet qui le soit. L'énoncé est d'abord annoncé « Cor. », le mot
étant ensuite barré et remplacé par « Th. 6 » : il a cessé en cours d'écriture
de le tenir pour un corollaire du précédent.}

C'est l'analogue, pour ces variétés abéliennes, du théorème de Lefschetz
faible : une section hyperplane voit la même chose que l'ambiant en dessous du
degré médian, et au degré médian elle en voit plus.

\pagerange{5}{5}
\section*{Équivalence numérique et $\tau$-équivalence (page 5)}

Le septième et dernier théorème quitte les jacobiennes. Deux cycles sont
\emph{numériquement} équivalents lorsqu'ils ont les mêmes nombres
d'intersection contre tout cycle complémentaire — c'est la relation la plus
grossière qu'on puisse raisonnablement écrire. Ils sont
\emph{$\tau$-équivalents} lorsqu'un multiple non nul de leur différence est
algébriquement équivalent à zéro — c'est l'équivalence algébrique, à la torsion
près. Le feuillet affirme que les deux coïncident.

Pour les \emph{diviseurs}, c'est un théorème : un diviseur numériquement
équivalent à zéro sur une variété projective a un multiple algébriquement
équivalent à zéro. C'est le résultat de Matsusaka (1957), et c'est ce que le
nom porté en marge du feuillet enregistre.

Le feuillet, lui, écrit l'énoncé pour les « cycles algébriques sur $X$
projective non singulière », sans restriction de codimension. Sous cette forme
il est faux.\footnote{Griffiths (1969) a montré que sur une hypersurface
quintique générale de $\mathbb{P}^{4}$ il existe des courbes homologiquement
équivalentes à zéro dont l'image par l'application d'Abel--Jacobi n'est pas de
torsion, donc dont aucun multiple n'est algébriquement équivalent à zéro.
Homologiquement équivalent à zéro entraîne numériquement équivalent à zéro : ces
courbes sont donc numériquement triviales sans être $\tau$-équivalentes à zéro,
et le théorème 7 tombe dès la codimension $2$. Le quotient qui mesure l'écart
s'appelle depuis le \emph{groupe de Griffiths}, et Clemens a montré plus tard
qu'il pouvait être infiniment engendré. L'énoncé est daté de 1964 et la
réfutation de 1969 : il n'y avait alors, sur cette question, aucun exemple dans
un sens ni dans l'autre.}

\pagerange{7}{8}
\section*{La lettre à Serre, 12 août 1964 (pages 7 et 8)}

\subsection*{Ce qu'elle annonce}

La lettre reprend le programme en prose et lui ajoute deux choses : la raison de
la borne du théorème 2, et une réduction complète du théorème 4 à un énoncé sur
les variétés abéliennes.

La raison de la borne est cohomologique. Le module de Tate $T_{\ell}(J^{i}(X))$
est, écrit-il, un quotient d'un sous-module de $H^{2i-1}(X,
\mathbb{Z}_{\ell}(i))$ — « et quand on sera plus savant, on saura prouver que
c'est même un sous-module dudit ». Le rang de ce module est au plus
$b_{2i-1}(X)$, et comme le module de Tate d'une variété abélienne de dimension
$g$ est de rang $2g$, on retrouve
$\dim J^{i}(X) \leq \tfrac{1}{2}b_{2i-1}(X)$.\footnote{Le souhait a été exaucé :
l'injectivité du module de Tate du représentant algébrique dans la cohomologie
$\ell$-adique du degré correspondant fait partie de ce que Murre a établi en
codimension $2$. La lettre ajoute qu'en caractéristique nulle on peut aussi bien
prendre la cohomologie de Hodge $H^{i}(X,\Omega^{i-1}_{X})$, ce qui est la
comparaison entre les deux constructions.} Une note marginale de la page 7,
d'une autre encre, ajoute une inclusion conjecturale des cycles dans la
cohomologie de Hodge ; sa lecture n'est pas assurée et elle n'est pas utilisée
ici.

\subsection*{La réduction de la positivité}

Voici la construction, qui est le cœur de la lettre. Soit $T$ une variété
de paramètres connexe, lisse, munie d'un point marqué $a$, et soit $z$ une
classe de cycles de codimension $m$ sur $T\times X$, nulle en $a$. On considère
sur $T\times T\times X$ les deux projections $p$ et $q$ sur $T\times X$ et la
projection $r$ sur $T\times T$, et on pose

\[
  D \;=\; r_{*}\!\left(p^{*}(z)\cdot q^{*}(z)\right),
\]

qui est une classe de diviseurs sur $T\times T$, c'est-à-dire une classe de
correspondance divisorielle. Elle est symétrique par construction. Comme elle
s'annule sur les fibres du point marqué, elle provient de la variété d'Albanese
$A = \mathrm{Alb}^{0}(T)$ — « si tu veux, tu peux supposer $T = A$ et $a$
l'origine », précise-t-il —, donc d'une classe de correspondance symétrique sur
$A\times A$. Une telle classe définit un morphisme $A \to \hat{A}$ vers la
variété duale ; soit $N$ son noyau. Sur le quotient $J = A/N$, la classe
descend en une correspondance symétrique \emph{non dégénérée}, et la question
est celle-ci, la seule de la lettre :

\begin{quote}
cette correspondance est-elle \emph{positive}, c'est-à-dire l'élément
correspondant de $\mathrm{NS}(J)$ est-il ample ?
\end{quote}

Une réponse affirmative donnerait le théorème 4 : la variété $J$ ainsi obtenue
est la jacobienne intermédiaire $J^{m}(X)$, et la classe est celle que son
autodualité définit.\footnote{Le passage de $A$ à $J=A/N$ est la même
construction que celle du théorème 1, écrite du côté de la correspondance
plutôt que du côté de la famille : $N$ y est ce que le théorème 1 note
$N(\xi)$. Sur $\mathbb{C}$, la positivité demandée est celle des relations de
Hodge--Riemann sur la cohomologie primitive de degré $2m-1$, et c'est elle qui
fait de la jacobienne intermédiaire de Weil une variété abélienne et non un
simple tore complexe. La marge de la page 7 rappelle, contre l'énoncé même, que
la classe doit être prise primitive.}

Il ajoute que la question suffit à révéler la méthode de construction générale
des $J^{i}$, et qu'on peut si l'on veut se borner au cas où $T$ est
$\mathrm{Pic}^{0}(Y)$ pour une sous-variété $Y$ de $X$ de codimension $m-1$,
les diviseurs de $Y$ algébriquement équivalents à zéro étant regardés comme
cycles de codimension $m$ de $X$ — ce qui est le mécanisme du théorème 5, vu
depuis l'autre bout.

\subsection*{Ce que la lettre dit de sa manière de travailler}

Deux phrases méritent d'être lues pour elles-mêmes. Sur la relation
d'équivalence qu'il ne sait pas décrire : « je ne sais pas bien identifier pour
l'instant, mais qui mérite certainement d'être appelée équivalence
d'Albanese ». Et sur l'ensemble : « je suis moralement sûr que ce sont les
bonnes, bien que je n'ai pas prouvé grand-chose dessus pour l'instant ; par
contre, j'ai des conjectures en pagaille ».

La certitude n'attend pas la démonstration, et le nom n'attend pas la
définition. Ce que le programme des sept théorèmes énonce n'est pas une liste de
résultats mais la description de l'objet qui devrait exister ; l'existence, il
la déduira de deux d'entre eux. Il demande enfin si « les spécialistes
abéliens » sauraient répondre à sa question — « peut-être Matsusaka ? Ou
toi-même ? » — et se justifie de la poser en disant qu'elle est « bien terre à
terre ».

\pagerange{2}{2}
\section*{Appendice I : la complétion d'un anneau gradué (page 2)}

Ce feuillet, barré de deux longues diagonales, appartient à un autre document
et n'a aucun rapport avec les jacobiennes. Il énonce ce qu'il advient de la
complétion sur un anneau gradué.

Soit $A = \bigoplus_{i\geq 0} A^{i}$ un anneau gradué et $\hat{A}$ son complété
pour la filtration par les degrés $\geq n$. Le feuillet écrit d'emblée

\[
  \hat{A} \;=\; \prod_{i \geq 0} A_{i},
\]

ce qui est l'identification à retenir : compléter un anneau gradué, c'est
remplacer la somme directe par le produit, c'est-à-dire autoriser les séries
infinies là où l'anneau gradué n'admettait que les
polynômes.\footnote{L'égalité demande que la filtration par les degrés
coïncide avec la filtration $\mathfrak{a}$-adique pour $\mathfrak{a} = A_{+}$,
donc que $A$ soit engendré en degré $1$ sur $A_{0}$. La page pose bien une
hypothèse à cet endroit, entre crochets, mais elle est en partie illisible : ce
qui s'en lit est « $A_{1}$ de type fini sur $A_{0}$ », ce qui est la moitié de
la condition. La lecture ne supplée pas l'autre moitié et signale seulement que
l'énoncé la requiert.}

Sous une hypothèse noethérienne, $\hat{A}$ est plat sur $A$, et l'on pose
$\hat{M} = M\otimes_{A}\hat{A}$ pour tout $A$-module $M$. Le feuillet établit
alors deux choses. D'abord, pour $M$ gradué de type fini, la même description
vaut :

\[
  \hat{M} \;=\; \prod_{i\geq 0} M_{i}.
\]

Ensuite, pour $M$ seulement de type fini, on passe par une présentation
$\cdots \to M''\to M\to M'\to 0$ avec $M''$ de type fini, on la complète — ce
que la platitude autorise —, et l'on obtient

\[
  \hat{M}'' \;=\; \prod_{i\geq 0} M''_{i} \;=\; \prod_{i\geq 0} M_{i},
  \qquad
  \hat{M}' \;\overset{\sim}{\longrightarrow}\; M' .
\]

Le second isomorphisme est ce qui rend l'argument concluant : la partie de $M$
qui n'était pas graduée ne voit pas la complétion.\footnote{Le raisonnement est
écrit très vite et sa prose ne se laisse lire qu'en partie ; la lecture ne
restitue que les formules, qui portent, et la charpente de l'argument, qui s'en
déduit. La dernière ligne du feuillet affirme que le foncteur de complétion est
exact et fidèle, ce qui est la conséquence de la platitude fidèle et non
seulement de la platitude.}

\pagerange{4}{4}
\section*{Appendice II : les axiomes d'une donnée radicielle (page 4)}

Ce second feuillet étranger, barré lui aussi, énumère deux structures
axiomatiques. La première est celle d'un système de racines réduit ; la seconde
est celle qu'on appelle aujourd'hui une \emph{donnée radicielle}.

\subsection*{Un système de racines, et l'intégralité des nombres de Cartan}

On se donne un couple $(P, \Delta)$ où $P$ est un groupe abélien libre de rang
fini et $\Delta$ une partie finie de $P$, soumis à trois conditions :
$\Delta$ engendre $P$ ; si $a \in \Delta$ alors $-a\in\Delta$ mais
$na\notin\Delta$ pour $n\neq \pm 1$ ; et pour tout $a\in\Delta$ il existe une
« symétrie » $S_{a}$ envoyant $a$ sur $-a$ et préservant
$\Delta$.\footnote{La deuxième condition est celle qui rend le système
\emph{réduit}. La première, qui demande que $\Delta$ engendre $P$, place
d'emblée dans le cas semi-simple : le formalisme de SGA 3 ne l'exige pas, et
c'est ce qui lui permet de couvrir les groupes réductifs quelconques.}

La note qui suit est le cœur du feuillet. Elle observe que les symétries
sont alors uniquement déterminées, parce qu'on peut les réaliser toutes à la
fois au moyen d'une seule forme quadratique invariante, écrite explicitement :

\[
  (x,y) \;=\; \sum_{a\in\Delta} \langle x, a\rangle\,\langle y, a\rangle
  \qquad \text{sur } (P')^{\mathbb{Q}} .
\]

De là, $S_{\alpha}x = x - \langle\gamma_{\alpha}, x\rangle\,\alpha$ pour un
$\gamma_{\alpha}\in (P')^{\mathbb{Q}}$ bien déterminé — c'est la coracine
$\alpha^{\vee}$ — et le feuillet affirme qu'en fait $\gamma_{\alpha}$ est
entier, c'est-à-dire que les nombres

\[
  c_{\alpha\beta} \;=\; \langle\gamma_{\alpha},\beta\rangle
\]

sont dans $\mathbb{Z}$ : ce sont les entiers de Cartan. La raison qu'il en donne
tient à ce que $\alpha$ soit indivisible dans $P$, du fait de pouvoir faire
partie d'un « système fondamental de racines ».\footnote{Ce dernier crochet est
en partie illisible et la lecture ne le reconstitue pas au-delà de sa forme
générale. L'ordre choisi est remarquable en lui-même : dans la présentation
usuelle l'intégralité des $\langle\alpha^{\vee},\beta\rangle$ est un
\emph{axiome} ; ici elle est une \emph{conséquence}, tirée de la forme
quadratique canonique et de l'indivisibilité des racines.}

\subsection*{La donnée radicielle}

La seconde structure est un quintuplet $(P, P', \Delta, \Delta', \varphi)$ où
$P$ et $P'$ sont deux groupes abéliens libres de rang fini en dualité,
$\Delta\subset P$, $\Delta'\subset P'$, et $\varphi$ une bijection de $\Delta$
sur $\Delta'$. Les conditions demandées sont que $\Delta$ engendre $P$ et
$\Delta'$ engendre $P'$, et que les deux familles de symétries se
correspondent : pour $\alpha\in\Delta$ d'image $\alpha' = \varphi(\alpha)$,

\[
  S_{\alpha} : x \longmapsto x - \langle \alpha', x\rangle\,\alpha
  \quad\text{laisse } \Delta \text{ stable},
  \qquad
  S'_{\alpha} : x' \longmapsto x' - \langle x', \alpha\rangle\,\alpha'
  \quad\text{laisse } \Delta' \text{ stable},
\]

la marge ajoutant la normalisation $\langle \alpha', \alpha\rangle = 2$.

C'est, à la condition d'engendrement près, la définition d'une donnée
radicielle.\footnote{La définition publiée est celle de l'exposé XXI de SGA 3,
rédigé par Demazure, où le quintuplet est noté $(M, M^{*}, R, R^{\vee})$ : $M$
le groupe des caractères d'un tore maximal, $M^{*}$ celui des cocaractères,
$R$ les racines et $R^{\vee}$ les coracines. La bijection $\varphi$ de ce
feuillet est $\alpha\mapsto\alpha^{\vee}$, et la normalisation
$\langle\alpha^{\vee},\alpha\rangle = 2$ y est l'un des deux axiomes. La
différence tient à la condition d'engendrement, que ce feuillet impose et que
SGA 3 lève : c'est elle qui sépare le cas semi-simple du cas réductif général,
et c'est la levée qui a fait la portée de la notion. La date du dossier — 1964
— est celle de la rédaction de SGA 3, et ce feuillet est du même côté du
tournant.}

\end{document}
